\documentclass[a4paper,12pt]{article}
\usepackage[italian]{babel}
\usepackage[utf8]{inputenc}
\usepackage[T1]{fontenc}
\usepackage{geometry}
\usepackage{setspace}
\usepackage{enumitem}
\usepackage{listings}
\usepackage{xcolor}
\geometry{margin=2.5cm}
\onehalfspacing

\definecolor{codegray}{rgb}{0.95,0.95,0.95}
\lstset{
    backgroundcolor=\color{codegray},
    basicstyle=\ttfamily\small,
    breaklines=true,
    frame=single,
    showstringspaces=false,
    tabsize=4,
    language=SQL
}

\title{Queries SQL}
\author{Prof. Fedeli Massimo}
\date{}

\begin{document}

\maketitle

\section*{Schema logico di riferimento}

\begin{flushleft}
STUDENTE(\underline{Id\_studente}, Cognome, Nome, Data\_nascita, Luogo\_nascita, CF, citta, CAP, via, n\_civico, Email\_istituzionale, Classe, Indirizzo\_studi) \\
AZIENDA(\underline{P\_iva}, Ragione\_sociale, Sede\_legale, Sede\_operativa, settore\_attivita, Email\_aziendale, Telefono\_ref) \\
TUTOR\_SCUOLA(\underline{Codice\_docente}, Cognome, Nome, Email\_istituzionale, Dipartimento\_app, N\_max\_studenti\_seguiti) \\
TUTOR\_AZIENDA(\underline{Codice\_dipendente}, Cognome, Nome, Email\_aziendale, P\_iva) \\
COMPETENZE(\underline{nome\_competenza}) \\
DISPONIBILITA(\underline{Codice\_id}, Data\_inizio, Data\_fine, N\_max\_studenti\_ospitabili, Descrizione\_attivita, Indirizzo\_studi\_consigliato) \\
AZIENDA\_DISPONIBILITA(\underline{P\_iva, Codice\_id}) \\
RICHIEDE(\underline{nome\_competenza, Codice\_id}) \\
ESPERIENZA(\underline{Id\_studente, P\_iva\_azienda\_ospitante, Codice\_docente, Codice\_dipendente, Codice\_id\_disponibilita}, Periodo\_effettivo, N\_ore\_previste, N\_ore\_svolte, Valutazione\_finale, esito)
\end{flushleft}


\section*{Quesiti}\begin{enumerate}[leftmargin=0.8cm,label=\textbf{\arabic*.}]
\item Elencare cognome, nome, classe ed email istituzionale di tutti gli studenti, ordinati per cognome e nome.
\item Visualizzare ragione sociale, settore di attività e sede operativa di tutte le aziende con sede operativa a Pesaro.
\item Mostrare i tutor scolastici che possono seguire almeno 11 studenti.
\item Elencare le disponibilità con Data fine successiva al 31 maggio 2026, mostrando codice, descrizione attività e numero massimo di studenti ospitabili.
\item Mostrare gli studenti della classe 5A nati dopo il 1 gennaio 2005.
\item Elencare per ogni esperienza lo studente, l'azienda ospitante e le ore previste.
\item Mostrare gli studenti che hanno svolto meno di 100 ore effettive, indicando anche azienda ospitante e ore svolte.
\item Elencare le aziende con la relativa disponibilità pubblicata, mostrando ragione sociale, descrizione attività e indirizzo di studi consigliato.
\item Visualizzare per ogni disponibilità le competenze richieste, mostrando codice disponibilità, descrizione attività e nome competenza.
\item Mostrare per ogni esperienza lo studente, il tutor scolastico e il tutor aziendale associati.
\item Calcolare quante aziende appartengono a ciascun settore di attività.
\item Calcolare il numero di studenti per ciascuna classe.
\item Per ogni azienda ospitante, calcolare il numero di esperienze svolte.
\item Determinare, per ogni esperienza, il deficit di ore rispetto a quelle previste, mostrando studente, ore previste, ore svolte e differenza.
\item Calcolare la media delle ore svolte per ciascun esito finale.
\item Trovare gli studenti che hanno svolto un'esperienza presso un'azienda del settore 'Software'.
\item Elencare le aziende che non appartengono ai settori 'IT' e 'Software'.
\item Mostrare le disponibilità il cui numero massimo di studenti ospitabili è maggiore della media complessiva.
\item Trovare gli studenti che hanno svolto un numero di ore superiore alla media delle ore svolte da tutti gli studenti.
\item Elencare i tutor scolastici che non risultano assegnati ad alcuna esperienza.
\item Determinare il settore di attività con il maggior numero di aziende.
\item Trovare l'azienda o le aziende che ospitano studenti con il numero massimo di ore svolte.
\item Per ogni settore di attività, calcolare la media delle ore svolte dagli studenti ospitati dalle aziende di quel settore.
\item Elencare le disponibilità associate ad aziende del settore 'AI' oppure 'Security', mostrando azienda, attività e competenza richiesta.
\item Trovare gli studenti che hanno svolto meno ore del corrispondente valore medio delle esperienze della loro stessa classe.\end{enumerate}
\end{document}
